\chapter{Conclusão e trabalhos futuros}
\label{cap:conclusao}
% --------------------------------------------------------------------------

O presente Trabalho de Conclusão de Curso estabeleceu, implementou e validou um arcabouço computacional autônomo e de alta fidelidade para o projeto, a aceleração preditiva neural e a otimização multiobjetivo de chassis estruturais primários para nanossatélites da classe CubeSat 1U. Ao integrar a tecnologia de metamateriais celulares com coeficiente de Poisson negativo (auxéticos reentrantes) à manufatura aditiva metálica por Fusão em Leito de Pó a Laser (L-PBF) em liga aeroespacial AlSi10Mg, a pesquisa demonstrou a viabilidade de superar o paradigma dos chassis maciços usinados convencionais, promovendo atenuação dinâmica passiva sem a introdução de massa parasitária.

\section{Síntese das conquistas e validação das hipóteses de pesquisa}
\label{sec:conquistas}

A condução sistemática da metodologia atendeu integralmente ao objetivo geral e aos nove objetivos específicos delineados no Capítulo~\ref{cap:objetivos}. Os resultados quantitativos obtidos na campanha computacional corroboraram as três hipóteses de pesquisa que nortearam a investigação:

\begin{alineas}
  \item \textbf{confirmação da Hipótese 1 (alívio de massa e preservação da rigidez dinâmica):} demonstrou-se que a substituição de paredes contínuas por matrizes celulares auxéticas de espessura de costela graduada promoveu uma redução de \textbf{\SI{62.49}{\percent}} na massa estrutural total do chassi 1U (de \SI{0.308}{\kilo\gram} no chassi monolítico usinado em Al 6061-T6 para apenas \textbf{\SI{0.116}{\kilo\gram}} no chassi ótimo proposto, poupando \SI{192.0}{\gram} do orçamento de lançamento). Simultaneamente, a hibridização entre o núcleo celular e os quatro trilhos-guia maciços de \rail{} elevou a primeira frequência natural fundamental para \textbf{\SI{579.62}{\hertz}} na simulação física de alta ordem (acréscimo de \SI{+13.20}{\percent} frente aos \SI{512.0}{\hertz} do chassi convencional), superando em quase seis vezes o limite mínimo de rigidez estabelecido pela norma NASA GSFC-STD-7000A ($f_1 \ge \SI{100.0}{\hertz}$);
  
  \item \textbf{confirmação da Hipótese 2 (atenuação vibratória passiva de alta performance):} o projeto de voo ótimo (\#28) exibiu coeficiente de Poisson efetivo acentuadamente negativo ($\nu_{\text{eff}} = \num{-13.81}$). Sob o espectro de vibração aleatória de qualificação espacial da NASA GEVS (\SI{14.1}{\Grms}, \SIrange{20}{2000}{\hertz}), a resposta elastodinâmica celular operou como um filtro mecânico passivo por antirressonância e desacoplamento de impedância acústica (e não por amortecimento viscoso arbitrário), reduzindo a transmissibilidade dinâmica de \num{0.850} para \textbf{\num{0.1997}} (atenuação espectral de \textbf{\SI{76.50}{\percent}} ou \textbf{\SI{12.6}{\decibel}}). Em consequência, a aceleração estocástica eficaz transmitida para a interface da carga útil reduziu-se acentuadamente de \SI{12.00}{\Grms} para \textbf{\SI{2.82}{\Grms}}, proporcionando um ambiente mecânico favorável para sensores ópticos, relógios atômicos e placas aviônicas sem recorrer a coxins viscoelásticos externos sujeitos a desgaseificação (\textit{outgassing});
  
  \item \textbf{confirmação da Hipótese 3 (aceleração evolutiva por metamodelo neural guiado por física):} o desenvolvimento do modelo substituto neural residual \code{CubeSatSurrogateResNet}, incorporando perdas físicas informadas pelas equações de escala de Gibson-Ashby e restrições de monotonicidade, atingiu coeficiente de determinação médio $\overline{R^2} = \num{0.9874}$ e erro percentual de apenas \textbf{\SI{+0.99}{\percent}} no pico de tensão estocástica $3\sigma$ e de \textbf{\SI{-0.008}{\percent}} na transmissibilidade dinâmica frente ao \textit{solver} Code\_Aster com elementos Gmsh Tet10. Essa precisão viabilizou a execução de \num{10000} avaliações do algoritmo genético NSGA-II em apenas \textbf{\SI{4.82}{\second}} ($\sim \SI{0.405}{\milli\second}$ por indivíduo), proporcionando um \textit{speedup} computacional superior a \num{200000} vezes ($\approx \num{207469}\times$) e poupando aproximadamente \SI{277.8}{\hour} ($\approx 11{,}6$ dias ininterruptos) de simulações contínuas em elementos finitos.
\end{alineas}

No tocante à integridade estrutural aeroespacial, o chassi ótimo satisfez estritamente todos os critérios de aeronavegabilidade e qualificação de voo:
\begin{alineas}
  \item \textbf{margem de segurança elástica estocástica ($MS_{\text{yield}} = \num{+0.102} > 0$):} a tensão combinada de von Mises a $3\sigma$ atingiu \SI{166.91}{\mega\pascal}, mantendo-se estritamente abaixo da tensão de projeto admissível com fator de qualificação ($\sigma_{\text{adm}} = \sigma_{\text{yield}} / 1{,}25 = \SI{184.0}{\mega\pascal}$ para a liga AlSi10Mg tratada termicamente);
  \item \textbf{vida em fadiga por vibração aleatória ($D = \num{0.0458} \ll \num{0.25}$):} o dano acumulado calculado pelo método das três bandas de Steinberg para \SI{120.0}{\second} de ensaio dinâmico por eixo foi inferior a um quinto do limite normativo de corte estabelecido pela norma NASA-HDBK-7005 ($D \le \num{0.25}$), garantindo fator de segurança superior a $5{,}4\times$ contra falha mecânica por fadiga vibracional de alto ciclo.
\end{alineas}

\section{Principais contribuições técnicas e metodológicas}
\label{sec:contribuicoes}

As contribuições originais proporcionadas por este trabalho à engenharia mecânica e aeroespacial compreendem:

\begin{alineas}
  \item \textbf{arquitetura híbrida de chassi CubeSat DfAM:} concepção inédita de um chassi estrutural primário 1U que harmoniza o núcleo aliviado de metamaterial auxético reentrante com quatro trilhos longitudinais sólidos maciços de \rail{}, garantindo estrita conformidade com as restrições cinemáticas e tribológicas do dispensador de lançamento P-POD ($Ra \le \SI{1.6}{\micro\metre}$, raios de adoçamento $R \ge \SI{1.0}{\milli\metre}$ e anodização dura superficial), ao mesmo tempo em que preserva canais contínuos de despoeiramento com diâmetro $d_{\text{drain}} \ge \SI{2.00}{\milli\metre}$ e vão livre superior a \SI{1.50}{\milli\metre} para total evacuação do pó não fundido;
  
  \item \textbf{paradigma \textit{Physics-Guided Deep Learning} para estruturas metamateriais:} introdução de uma formulação de perda composta ($\mathcal{L}_{\text{total}} = \mathcal{L}_{\text{MSE}} + \lambda_{\text{phys}}\mathcal{L}_{\text{GA}} + \lambda_{\text{mono}}\mathcal{L}_{\text{mono}}$) na rede neural residual profunda, forçando o metamodelo a respeitar a física assintótica dos sólidos celulares de Gibson-Ashby e mitigando comportamentos espúrios nas regiões limítrofes do espaço de projeto;
  
  \item \textbf{orquestração computacional determinística via grafos de fluxo de dados:} estruturação do fluxo completo de projeto, triagem e simulação em um grafo de estados determinístico fortemente tipado (\code{CubeSatState}) via LangGraph (\code{StateGraph}) e compatível com ferramentas abertas como Prefect. O mecanismo implementa validação antecipada de viabilidade DfAM e NASA GEVS com roteamento \textit{fail-fast}, assegurando rastreabilidade total de decisões e estabelecendo uma esteira reprodutível e escalável para engenharia de sistemas espaciais;
  
  \item \textbf{esteira computacional de qualificação aeroespacial 100\% \textit{open source}:} desenvolvimento e homologação de uma suíte computacional totalmente aberta, integrando CadQuery (parametrização B-Rep e portões DfAM), Gmsh (malhas de segunda ordem Tet10), Code\_Aster e CalculiX (dinâmica modal e resposta espectral PSD), análise de dispersão de Bloch-Floquet (comprovação de \textit{bandgaps}) e OpenRadioss (dinâmica transiente de choque de separação e espectro SRS), eliminando a dependência de licenças comerciais e viabilizando execução sem custos em clusters HPC.
\end{alineas}

\section{Limitações da pesquisa e recomendações para trabalhos futuros}
\label{sec:limitacoes}

A despeito do elevado nível de maturidade computacional e consistência física dos modelos apresentados, reconhecem-se limitações metodológicas intrínsecas à abordagem numérica e ao uso de materiais análogos (nível de maturidade tecnológica TRL~3/4). Para a progressão da tecnologia rumo à qualificação de voo orbital (TRL~5 a 7), recomendam-se os seguintes desdobramentos:

\begin{alineas}
  \item \textbf{ensaio dinâmico em mesa vibratória e correlação modal (GVT):} fabricação de protótipos físicos em liga de alumínio AlSi10Mg por L-PBF com alívio térmico de tensões e condução de ensaios de Teste de Vibração Terrestre (GVT - \textit{Ground Vibration Test}) e de mesa vibratória triaxial sob o perfil PSD da norma NASA GSFC-STD-7000A (\SI{14.1}{\Grms}, \SI{120}{\second}/eixo), superando a similaridade incompleta associada à prototipagem em PLA e calibrando experimentalmente as curvas de transmissibilidade;
  
  \item \textbf{ensaios experimentais de choque e separação do P-POD:} realização de ensaios de disparo em dispensador de teste instrumentado com acelerômetros piezoelétricos na carga útil e nos trilhos, confrontando a aceleração de pico de desaceleração e o espectro SRS medidos frente às predições do solver explícito OpenRadioss;
  
  \item \textbf{escalabilidade computacional massiva da esteira aberta em clusters HPC:} implantação dos fluxos paralelos da esteira \textit{open source} (CadQuery, Gmsh, Code\_Aster, OpenRadioss) em ambiente distribuído via orquestradores de contêineres e filas Slurm/MPI, expandindo a exploração DoE para dezenas de milhares de topologias celulares;
  
  \item \textbf{qualificação termoelástica em câmara de termovácuo (TVAC):} simulação numérica acoplada termo-mecânica e ensaios experimentais sob os ciclos térmicos severos da órbita baixa terrestre (LEO), variando entre \SI{-40}{\celsius} e \SI{+85}{\celsius} em ambiente de vácuo térmico ($p \le \SI{e-6}{\torr}$), investigando a dissipação térmica anisotrópica do núcleo auxético;
  
  \item \textbf{escalabilidade para plataformas multimodulares (CubeSat 3U, 6U e 12U):} expansão da formulação geométrica para formatos espaciais de maior envergadura, introduzindo Gradação Funcional de Densidade tridimensional (FGM) contínua otimizada topologicamente para suportar cargas concentradas provenientes de apêndices mecânicos complexos.
\end{alineas}
